\textbf{Phần 1. Trắc nghiệm lựa chọn}
\setcounterpageref{ex}{0}

\begin{ex} %[2D14V69]
Tính: $48 + 37$.
\choice
{$75$}
{\True $85$}
{$86$}
{$95$}
\loigiai{
$48 + 37 = 85$.
\textbf{Đáp án của câu hỏi là B}
}
\end{ex}

\begin{ex} %[2D14V69]
Tính: $92 - 46$.
\choice
{$56$}
{$48$}
{$36$}
{\True $46$}
\loigiai{
$92 - 46 = 46$.
\textbf{Đáp án của câu hỏi là D}
}
\end{ex}

\begin{ex} %[2D14V69]
Tìm $x$ biết: $x + 29 = 75$.
\choice
{\True $46$}
{$56$}
{$44$}
{$104$}
\loigiai{
$x = 75 - 29 = 46$.
\textbf{Đáp án của câu hỏi là A}
}
\end{ex}

\begin{ex} %[2D14V69]
Một lớp học có $35$ bạn học sinh, trong đó có $18$ bạn nữ. Lớp học có bao nhiêu bạn nam?
\choice
{$16$ bạn}
{$19$ bạn}
{\True $17$ bạn}
{$27$ bạn}
\loigiai{
$35 - 18 = 17$ (bạn).
\textbf{Đáp án của câu hỏi là C}
}
\end{ex}

\begin{ex} %[2D14V69]
Tổng của hai số là $80$, số hạng thứ nhất là $34$. Số hạng thứ hai là
\choice
{$56$}
{\True $46$}
{$44$}
{$36$}
\loigiai{
$80 - 34 = 46$.
\textbf{Đáp án của câu hỏi là B}
}
\end{ex}

\begin{ex} %[2D14V69]
Mai hái được $28$ bông hoa, Lan hái được nhiều hơn Mai $14$ bông hoa. Cả hai bạn hái được bao nhiêu bông hoa?
\choice
{$42$ bông}
{$68$ bông}
{$72$ bông}
{\True $70$ bông}
\loigiai{
Lan hái: $28 + 14 = 42$ bông. Cả hai bạn: $28 + 42 = 70$ bông hoa.
\textbf{Đáp án của câu hỏi là D}
}
\end{ex}

\begin{ex} %[2D14V69]
Hiệu của hai số là $25$. Nếu số bị trừ tăng $7$ đơn vị và số trừ giữ nguyên thì hiệu mới là
\choice
{\True $32$}
{$18$}
{$30$}
{$35$}
\loigiai{
Số bị trừ tăng thì hiệu tăng: $25 + 7 = 32$.
\textbf{Đáp án của câu hỏi là A}
}
\end{ex}

\begin{ex} %[2D14V69]
Giá trị của biểu thức $100 - 45 - 25$ là
\choice
{$40$}
{$20$}
{\True $30$}
{$35$}
\loigiai{
$100 - (45 + 25) = 100 - 70 = 30$.
\textbf{Đáp án của câu hỏi là C}
}
\end{ex}

\begin{ex} %[2D14V69]
Điền dấu thích hợp: $56 + 28 \dots 94 - 10$.
\choice
{$>$}
{\True $=$}
{$<$}
{Không so sánh được}
\loigiai{
$56 + 28 = 84$; $94 - 10 = 84$. Hai vế bằng nhau.
\textbf{Đáp án của câu hỏi là B}
}
\end{ex}

\begin{ex} %[2D14V69]
Bác Ba có $62$ con gà, bác bán đi $25$ con gà. Sau đó bác mua thêm $18$ con gà nữa. Bác Ba có tất cả bao nhiêu con gà?
\choice
{$45$ con}
{$65$ con}
{$57$ con}
{\True $55$ con}
\loigiai{
$62 - 25 + 18 = 37 + 18 = 55$ (con).
\textbf{Đáp án của câu hỏi là D}
}
\end{ex}

\begin{ex} %[2D14V69]
Tìm số tự nhiên có hai chữ số biết chữ số hàng chục hơn hàng đơn vị 4 đơn vị và tổng hai chữ số bằng 12.
\choice
{\True $84$}
{$73$}
{$95$}
{$62$}
\loigiai{
Hàng chục: $(12 + 4) : 2 = 8$. Hàng đơn vị: $8 - 4 = 4$. Số đó là 84.
\textbf{Đáp án của câu hỏi là A}
}
\end{ex}

\begin{ex} %[2D14V69]
Một cuộn dây dài $80\text{ m}$. Người ta cắt đi $35\text{ m}$ lần đầu và $20\text{ m}$ lần sau. Cuộn dây còn lại bao nhiêu mét?
\choice
{$35\text{ m}$}
{$15\text{ m}$}
{\True $25\text{ m}$}
{$30\text{ m}$}
\loigiai{
$80 - 35 - 20 = 25$ (m).
\textbf{Đáp án của câu hỏi là C}
}
\end{ex}

\textbf{Phần 2. Trắc nghiệm đúng sai}
\setcounterpageref{ex}{0}

\begin{ex} %[2D14V69]
Xét tính đúng sai của các phép toán cộng trừ trong phạm vi 100:
\choiceTF
{\True 37 + 63 = 100.}
{81 - 29 = 62.}
{\True Trong phép trừ có nhớ phạm vi 100, ta mượn 1 chục ở hàng chục của số bị trừ.}
{\True Nếu a + b = 90 thì b = 90 - a.}
\loigiai{
\begin{itemize}
    \item Ý a) Đúng: Đúng, tổng là 100.
    \item Ý b) Sai: 81 - 29 = 52, không phải 62.
    \item Ý c) Đúng: Đúng quy tắc.
    \item Ý d) Đúng: Đúng tính chất tìm số hạng.
\end{itemize}
\begin{itemize}
    \item a) \textbf{Đúng}.
    \item b) \textbf{Sai}.
    \item c) \textbf{Đúng}.
    \item d) \textbf{Đúng}.
\end{itemize}
}
\end{ex}

\begin{ex} %[2D14V69]
Một cửa hàng có 75 kg gạo, buổi sáng bán 28 kg, buổi chiều bán 32 kg:
\choiceTF
{\True Cả hai buổi bán được 60 kg gạo.}
{\True Cửa hàng còn lại 15 kg gạo.}
{Số gạo bán buổi chiều ít hơn buổi sáng.}
{\True Buổi chiều bán nhiều hơn buổi sáng 4 kg.}
\loigiai{
\begin{itemize}
    \item Ý a) Đúng: 28 + 32 = 60 kg. Đúng.
    \item Ý b) Đúng: 75 - 60 = 15 kg. Đúng.
    \item Ý c) Sai: 32 > 28 nên bán nhiều hơn. Sai.
    \item Ý d) Đúng: 32 - 28 = 4 kg. Đúng.
\end{itemize}
\begin{itemize}
    \item a) \textbf{Đúng}.
    \item b) \textbf{Đúng}.
    \item c) \textbf{Sai}.
    \item d) \textbf{Đúng}.
\end{itemize}
}
\end{ex}

\textbf{Phần 3. Trả lời ngắn}
\setcounterpageref{ex}{0}

\begin{ex} %[2D14V69]
Tính: $64 + 28 - 45$.
\shortans{47}
\loigiai{
\textbf{Đáp số:} $47$\\
$92 - 45 = 47$.
}
\end{ex}

\begin{ex} %[2D14V69]
Tìm $x$ biết $100 - x = 38$.
\shortans{62}
\loigiai{
\textbf{Đáp số:} $62$\\
$x = 100 - 38 = 62$.
}
\end{ex}

\begin{ex} %[2D14V69]
Nam có 45 viên bi, Việt có ít hơn Nam 12 viên bi. Cả hai bạn có bao nhiêu viên bi?
\shortans{78}
\loigiai{
\textbf{Đáp số:} $78$\\
Việt có: $45 - 12 = 33$. Cả hai bạn: $45 + 33 = 78$ viên.
}
\end{ex}

\begin{ex} %[2D14V69]
Tổng của hai số là 72, một số là 39. Tìm số còn lại.
\shortans{33}
\loigiai{
\textbf{Đáp số:} $33$\\
$72 - 39 = 33$.
}
\end{ex}

\begin{ex} %[2D14V69]
Một bến xe có 54 xe khách, có 18 xe rời bến và 9 xe vào bến. Bến xe có bao nhiêu xe?
\shortans{45}
\loigiai{
\textbf{Đáp số:} $45$\\
$54 - 18 + 9 = 45$ xe.
}
\end{ex}

\begin{ex} %[2D14V69]
Tính hiệu của số chẵn lớn nhất có hai chữ số và số 45.
\shortans{53}
\loigiai{
\textbf{Đáp số:} $53$\\
$98 - 45 = 53$.
}
\end{ex}
